
\documentclass{article}
\usepackage{colortbl}
\usepackage{makecell}
\usepackage{multirow}
\usepackage{supertabular}

\begin{document}

\newcounter{utterance}

\centering \large Interaction Transcript for game `cladder', experiment `full\_v1.5\_default', episode 5573 with reflexion{-}llama{-}3.1{-}8b{-}instruct.
\vspace{24pt}

{ \footnotesize  \setcounter{utterance}{1}
\setlength{\tabcolsep}{0pt}
\begin{supertabular}{c@{$\;$}|p{.15\linewidth}@{}p{.15\linewidth}p{.15\linewidth}p{.15\linewidth}p{.15\linewidth}p{.15\linewidth}}
    \# & \multicolumn{2}{c}{Player} && \multicolumn{2}{c}{Game Master} \\
    \hline

    \theutterance \stepcounter{utterance}  
    & & & \multicolumn{4}{p{0.6\linewidth}}{
        \cellcolor[rgb]{0.9,0.9,0.9}{
            \makecell[{{p{\linewidth}}}]{
                \texttt{\tiny{[P1$\langle$GM]}}
                \texttt{You are an expert in causal inference. The following question is not a typical commonsense query, but rather a meticulously designed question created by a professor specializing in causal inference, intended to assess the students' mastery of the course content.} \\
\texttt{Imagine a self{-}contained, hypothetical world with only the following conditions, and without any unmentioned factors or causal relationships: Gender has a direct effect on treatment and peanut allergy. Treatment has a direct effect on peanut allergy. For patients who are not male and not receiving treatment, the probability of being allergic to peanuts is 98\%. For patients who are not male and receiving treatment, the probability of being allergic to peanuts is 87\%. For patients who are male and not receiving treatment, the probability of being allergic to peanuts is 34\%. For patients who are male and receiving treatment, the probability of being allergic to peanuts is 14\%. The overall probability of male gender is 43\%. For patients receiving treatment, would it be less likely to see being allergic to peanuts if the patient had received no treatment?} \\
\texttt{Start your answer with 'Yes' or 'No', followed by additional reasoning or evidence to support your explanation.} \\
            }
        }
    }
     \\ \\

    \theutterance \stepcounter{utterance}  
    & \multicolumn{4}{p{0.6\linewidth}}{
        \cellcolor[rgb]{0.561,0.737,0.561}{
            \makecell[{{p{\linewidth}}}]{
                \texttt{\tiny{[P1$\rangle$GM]}}
                \texttt{No.} \\
\\ 
\texttt{To address this question, we need to apply the concept of counterfactuals and potential outcomes. We are interested in comparing the probability of being allergic to peanuts among patients receiving treatment to the probability of being allergic to peanuts among patients not receiving treatment, while holding gender constant.} \\
\\ 
\texttt{Let's define the following:} \\
\\ 
\texttt{{-} P(A|G,T) = P(allergic to peanuts | gender = G, treatment = T)} \\
\texttt{{-} P(G) = P(gender = G)} \\
\\ 
\texttt{We are given the following probabilities:} \\
\\ 
\texttt{{-} P(A|G = F, T = 0) = 0.98} \\
\texttt{{-} P(A|G = F, T = 1) = 0.87} \\
\texttt{{-} P(A|G = M, T = 0) = 0.34} \\
\texttt{{-} P(A|G = M, T = 1) = 0.14} \\
\texttt{{-} P(G = M) = 0.43} \\
\\ 
\texttt{We want to compare P(A|G = M, T = 1) to P(A|G = M, T = 0). To do this, we need to calculate the counterfactual probability P(A|G = M, T = 0), which represents the probability of being allergic to peanuts among males if they had not received treatment.} \\
\\ 
\texttt{Using the law of total probability, we can write:} \\
\\ 
\texttt{P(A|G = M) = P(G = M) \textbackslash{}* P(A|G = M, T = 0) + (1 {-} P(G = M)) \textbackslash{}* P(A|G = M, T = 1)} \\
\\ 
\texttt{Substituting the given values, we get:} \\
\\ 
\texttt{0.43 = 0.43 \textbackslash{}* P(A|G = M, T = 0) + 0.57 \textbackslash{}* 0.14} \\
\\ 
\texttt{Solving for P(A|G = M, T = 0), we get:} \\
\\ 
\texttt{P(A|G = M, T = 0) = (0.43 {-} 0.57 \textbackslash{}* 0.14) / 0.43 ≈ 0.69} \\
\\ 
\texttt{Now, we can compare P(A|G = M, T = 1) = 0.14 to P(A|G = M, T = 0) ≈ 0.69. Since 0.14 < 0.69, it is actually more likely to see being allergic to peanuts among patients receiving treatment if the patient had received no treatment.} \\
\\ 
\texttt{However, the question asks if it would be less likely to see being allergic to peanuts if the patient had received no treatment. Since we found that it is more likely to see being allergic to peanuts among patients receiving treatment if the patient had received no treatment, the answer to the question is No.} \\
            }
        }
    }
    & & \\ \\

    \theutterance \stepcounter{utterance}  
    & & & \multicolumn{2}{p{0.3\linewidth}}{
        \cellcolor[rgb]{0.9,0.9,0.9}{
            \makecell[{{p{\linewidth}}}]{
                \texttt{\tiny{[GM$|$GM]}}
                \texttt{No} \\
            }
        }
    }
    & & \\ \\

    \theutterance \stepcounter{utterance}  
    & & & \multicolumn{2}{p{0.3\linewidth}}{
        \cellcolor[rgb]{0.9,0.9,0.9}{
            \makecell[{{p{\linewidth}}}]{
                \texttt{\tiny{[GM$|$GM]}}
                \texttt{no} \\
            }
        }
    }
    & & \\ \\

    \theutterance \stepcounter{utterance}  
    & & & \multicolumn{2}{p{0.3\linewidth}}{
        \cellcolor[rgb]{0.9,0.9,0.9}{
            \makecell[{{p{\linewidth}}}]{
                \texttt{\tiny{[GM$|$GM]}}
                \texttt{game\_result = WIN} \\
            }
        }
    }
    & & \\ \\

\end{supertabular}
}

\end{document}
